% =====================================================================
%  25.4  Beyond Paths -- virtual children, subtree aggregates, and the
%  connectivity question (LCT vs Euler-tour tree vs DSU).
%  Exports: sec:c23-beyond  def:c23-virtual  def:c23-map
%           thm:c23-vsum  lst:c23-subtree  fig:c23-virtual
%  WORKED examples only -- NO solvedproblem (in-section solved stays 2).
% =====================================================================
\section{Beyond Paths: Virtual Children and Connectivity}\label{sec:c23-beyond}

\begin{previouslybox}[Paths were easy because a splay only sees the solid chain]
\(\secref{c23-makeroot}\)'s \code{query} read a path and was done, because a
preferred path \emph{is} the splay subtree -- the off-path branches simply did not
matter. The moment you ask for a \emph{subtree} aggregate on a moving tree, those
off-path branches matter enormously: an \code{access} is constantly promoting one
child to solid and demoting the others to \emph{virtual}, and a node's splay
children are only ever its preferred neighbours. This section shows the fix -- a
separate accumulator \(\mathit{vir}[\cdot]\) for virtual children, kept current on
every access -- and then steps back to ask, honestly, when you should not be
reaching for an LCT at all.
\end{previouslybox}

\begin{topicintro}
Two limits define the LCT's real job description. First, a subtree aggregate: a
node's aux splay holds only its \emph{preferred} child, so the other represented
children are ``virtual'' and invisible to \code{pull}; you fix this by storing their
combined value in \(\mathit{vir}[x]\) and moving a child between the solid and
virtual sets during \code{access} (\(\thmref{c23-vsum}\)). Second, connectivity: if the
query only asks ``are \(u,v\) in the same tree?'' -- never a path value -- an
Euler-tour tree or a (rollback) DSU is smaller, faster, and far less bug-prone.
Knowing both limits is what separates an LCT that ships from an LCT that TLEs on a
\(10^5\) connectivity problem it was never meant to solve.

\textbf{Prerequisites.} \(\secref{c23-makeroot}\) -- \code{access}/\code{makeroot}
and a \code{pull} monoid. Chapter~5 (\chref{c24}) -- the Euler tour that an
Euler-tour tree keeps in a balanced BST. Chapter~17 (\chref{c28}) -- DSU, and
Chapter~24 (\chref{c18}) -- the offline rollback / divide-and-conquer alternative.

\textbf{After this section you can}
\begin{itemize}[nosep]
  \item keep a virtual-child accumulator \(\mathit{vir}\) correct across accesses; \hfill\lv{5}
  \item answer a subtree-sum on a link/cut tree; \hfill\lv{5}
  \item decide LCT vs Euler-tour tree vs DSU from the query's shape. \hfill\lv{5}
\end{itemize}
\end{topicintro}

\begin{definition}{Virtual children and the \(\mathit{vir}\) accumulator}{c23-virtual}
Call a represented child of \(x\) \term{solid} if it lies on \(x\)'s current
preferred path (so it is one of \(x\)'s two aux children); every other represented
child is \term{virtual}. Define \(\mathit{vir}[x]\) as the fold (here a sum) over the
\emph{subtree-aggregate} of each virtual child. A node's full subtree value is then
\(\mathit{sub}[x]=\mathit{val}[x]+\mathit{vir}[x]+\mathit{sub}[\text{solid children}]\),
and the path fold \(\mathit{sum}[x]\) of \(\secref{c23-makeroot}\) is the special case
that omits \(\mathit{vir}\).
\end{definition}

\begin{theorem}{Access moves children between the solid and virtual sets}{c23-vsum}
In \code{access}, when at node \(x\) we set the right child from an old value \(r\) to
the new preferred child \(y\), \(r\) becomes virtual and \(y\) stops being virtual.
Maintaining
\[
  \mathit{vir}[x]\;\mathrel{-}=\;\mathit{sub}[y]\quad(\text{if }y\ne 0),\qquad
  \mathit{vir}[x]\;\mathrel{+}=\;\mathit{sub}[r]\quad(\text{if }r\ne 0)
\]
at that instant, and folding \(\mathit{vir}\) into \(\mathit{sub}\) (not into the pure
path \(\mathit{sum}\)) keeps every subtree aggregate exact after any access.
\end{theorem}
\begin{proofbox}
A subtree aggregate must count every represented child's subtree exactly once. At
any moment the represented children of \(x\) partition into the (at most one) solid
child on the preferred path and the rest, which are virtual. Before the access step,
\(r\) was the solid (right/deeper) child and \(y\) was virtual; after, \(y\) is solid
and \(r\) is virtual. Removing \(\mathit{sub}[y]\) from \(\mathit{vir}[x]\) (it is now
counted via the solid splay path instead) and adding \(\mathit{sub}[r]\) (it is no
longer on the path) preserves the invariant ``\(\mathit{vir}[x]\) is the sum over
precisely the virtual children's subtrees.'' Since \(\mathit{sub}[x]\) is defined as
\(\mathit{val}+\mathit{vir}+\) solid-subtrees, and every rotation of the aux splay only
reshuffles solid children (whose \(\mathit{sub}\) values are unchanged), \(\mathit{sub}\)
stays exact. \(\Box\)
\end{proofbox}

\begin{figure}[H]
\centering
\begin{tikzpicture}[diagram, scale=0.82, transform shape,
   snode/.style={circle, draw=primaryblue!70, fill=primaryblue!8, line width=0.7pt,
                 minimum size=6mm, font=\scriptsize\bfseries, inner sep=1pt},
   virt/.style={circle, draw=successgreen!75, fill=successgreen!10, line width=1pt,
              minimum size=6.4mm, font=\footnotesize\bfseries, inner sep=1pt},
   pref/.style={circle, draw=difforange, fill=difforange!16, line width=1pt,
              minimum size=6.4mm, font=\footnotesize\bfseries, inner sep=1pt}]
  \node[pref] (n2) at (0,0) {2};
  \node[pref] (n5) at (0,-1.5) {5};
  \node[virt] (n4) at (-2.2,-1.5) {4};
  \node[virt] (n6) at (2.2,-1.5) {6};
  \node[pref] (n7) at (0,-3.0) {7};
  \draw[difforange, line width=1.6pt] (n2)--(n5)--(n7);
  \draw[successgreen, line width=1.3pt, densely dashed] (n2)--(n4);
  \draw[successgreen, line width=1.3pt, densely dashed] (n5)--(n6);
  \node[annot, font=\scriptsize, text=successgreen, anchor=east] at (-2.4,-0.7) {virtual: in \(\mathit{vir}[2]\)};
  \node[annot, font=\scriptsize, text=successgreen, anchor=west] at (2.4,-0.7) {virtual: in \(\mathit{vir}[5]\)};
\end{tikzpicture}
\caption{On the running forest \(R\) (\(\figref{c23-tree}\)) after \code{access(7)},
the solid path \(2\!-\!5\!-\!7\) (orange) lives in the aux splay; children \(4\) (of
\(2\)) and \(6\) (of \(5\)) are \emph{virtual} (dashed green) and are excluded from
the path fold but included in \(\mathit{vir}[2]\), \(\mathit{vir}[5]\) so that a
\emph{subtree} query still counts them. When a later access demotes \(5\)'s solid
child and promotes \(6\), the two \(\mathit{sub}\) values trade between the splay and
\(\mathit{vir}[5]\) exactly as \(\thmref{c23-vsum}\) prescribes.}
\label{fig:c23-virtual}
\end{figure}

\begin{lstlisting}[caption={Subtree-sum-capable LCT: path fold sum[] (solid only) and subtree fold sub[] (solid + virtual) kept side by side. access() now updates vir[] as children swap between solid and virtual. ASCII only}, label={lst:c23-subtree}]
long long val[MAXN], sum[MAXN], vir[MAXN], sub[MAXN];
void pull(int x){
    sum[x] = val[x] + sum[ch[x][0]] + sum[ch[x][1]];        // PATH fold (no vir)
    sub[x] = val[x] + vir[x] + sub[ch[x][0]] + sub[ch[x][1]];// SUBTREE fold (+ vir)
}
void access(int x){
    for(int last=0; x; last=x, x=fa[x]){
        splay(x);
        int r = ch[x][1];                    // old solid (deeper) child -> becomes virtual
        if(r) vir[x] += sub[r];
        if(last) vir[x] -= sub[last];        // new solid child was virtual until now
        ch[x][1] = last;
        pull(x);
    }
}
long long subtree_sum(int x){ access(x); return sub[x]; }   // x's aux root holds all
\end{lstlisting}

\begin{insightbox}[Why the path fold ignores \(\mathit{vir}\) but the subtree fold needs it]
A path \(u\leadsto v\) is a single preferred chain -- the vertices on it are exactly
the splay subtree, and their off-path branches are NOT part of the path, so
\(\mathit{sum}\) must exclude \(\mathit{vir}\). A subtree, by contrast, is ``this
node and everything below it,'' which is precisely the branches \(\mathit{vir}\)
holds. Keeping the two folds as sibling fields in \code{pull} -- never mixing them --
is why the same structure can answer a path query and a subtree query without a
rebuild; the Summit will use only the path fold, but knowing the sibling makes that
choice deliberate rather than lucky.
\end{insightbox}

\subsection*{The connectivity question, and the LCT-vs-Euler-tour pick}

Not every dynamic-forest problem needs a path aggregate. If the only query is
``are \(u,v\) connected?'' (or ``is this graph still a forest after adding edge
\((u,v)\)?''), an LCT is overkill -- you pay splay constants for a boolean. An
Euler-tour tree stores the tour in a balanced BST and splits/merges on
\code{link}/\code{cut}, answering connectivity by testing whether two nodes lie in
one tour; a rollback DSU does it offline in near-linear time with a fraction of the
code. The decision is the pick table below.

\begin{definition}{Which dynamic-forest tool? (pick table)}{c23-map}
\emph{Path value needed online} (sum/min/max over the route) \(\Rightarrow\) full LCT
with \(\mathit{sum}\) (\(\secref{c23-makeroot}\)). \emph{Subtree value needed online}
\(\Rightarrow\) LCT with the \(\mathit{vir}\) sibling (\(\thmref{c23-vsum}\)). \emph{Only
connectivity / component identity, online} \(\Rightarrow\) Euler-tour tree (simpler
than LCT, \(\chref{c24}\)). \emph{Only connectivity, all edits known offline}
\(\Rightarrow\) rollback DSU or segment-tree-on-time (\(\chref{c28}\), engine
\(\chref{c18}\)). \emph{Fixed tree, path/subtree ranges} \(\Rightarrow\) not dynamic at
all: HLD (\(\chref{c16}\)) or Euler+Fenwick (\(\chref{c24}\)) -- see
\(\secref{c23-boundary}\).
\end{definition}

\begin{warningbox}[Subtree aggregates are where LCTs go to die]
The \(\mathit{vir}\) update in \(\coderef{c23-subtree}\) is easy to get subtly wrong:
forgetting to \emph{subtract} \(\mathit{sub}[\mathit{last}]\) (or to add
\(\mathit{sub}[r]\)) leaves a stale virtual sum that only shows up on subtree queries
whose access path demoted a child -- a rare, order-dependent WA. If your problem can
be solved by an Euler-tour tree or offline DSU, take that exit; the \(\mathit{vir}\)
discipline is the single hardest part of a production LCT and is exactly what
\(\secref{c23-boundary}\) tells you to avoid when a simpler tool fits.
\end{warningbox}

\begin{hardeningbox}
\begin{itemize}[nosep,leftmargin=1.4em]
  \item \textbf{Two folds, two uses.} Never read \(\mathit{sub}\) for a path query or
        \(\mathit{sum}\) for a subtree query; \(\mathit{sum}\) omits \(\mathit{vir}\),
        \(\mathit{sub}\) includes it. Mixing them silently mis-counts every branch.
  \item \textbf{Order inside access.} Compute \(\mathit{sub}[r]\) and
        \(\mathit{sub}[\mathit{last}]\) \emph{before} re-pointing \(\mathit{ch}[x][1]\),
        then \code{pull(x)} after -- otherwise the subtracted value is already stale.
  \item \textbf{Rooted vs rooted-after-evert.} \(\mathit{vir}\) is defined against the
        \emph{current} represented root; after \code{makeroot} the parent direction
        flips and every \(\mathit{vir}\) along the everted path changes meaning -- an
        evert plus a subtree query is legal but must recompute, not reuse cached
        \(\mathit{sub}\) from before the flip.
\end{itemize}
\end{hardeningbox}

\begin{recallbox}[25.4 in one breath]
\emph{A splay only sees solid (preferred) children; the virtual children live in
\(\mathit{vir}[x]\). \code{access} swaps a child between the two sets, so update
\(\mathit{vir}\) by moving \(\mathit{sub}[\cdot]\) values in and out
(\(\thmref{c23-vsum}\)); the subtree fold adds \(\mathit{vir}\), the path fold omits it.
And before any of this: is a path/subtree value actually required? If only
connectivity is, an Euler-tour tree or rollback DSU wins (\(\defref{c23-map}\)); if
the tree never changes, it was never an LCT problem.}
\selftest{
  \item Why does the path fold exclude \(\mathit{vir}\) but the subtree fold include it? \emph{(A path is one solid chain; a subtree includes the virtual branches, Def.~\defref{c23-virtual}.)}
  \item During \code{access}, what happens to the old right child \(r\) and the new child \(\mathit{last}\)? \emph{(\(r\) becomes virtual (+\(\mathit{sub}[r]\)), \(\mathit{last}\) becomes solid (\(-\mathit{sub}[\mathit{last}]\)), \(\thmref{c23-vsum}\).)}
  \item Name the two easier tools for connectivity-only dynamic forests. \emph{(Euler-tour tree; rollback DSU / offline D\&C, \(\defref{c23-map}\).)}
}
\end{recallbox}
